\part{Das Handwerk}

\chapter{Werkzeuge}\label{kap:werkzeuge}

\begin{merksatz}
\boxtitel{Was mitgeliefert wird}
\begin{center}
{\small
\renewcommand{\arraystretch}{1.5}
\begin{tabular}{@{}P{2.6cm}P{3.0cm}P{2.4cm}@{}}
\toprule
\textbf{Werkzeug} & \textbf{Leistet} & \textbf{Aufwand}\\
\midrule
\code{gofmt} & \textbf{Formatierung} & \textcolor{gruen}{keiner}\\
\code{go vet} & \textbf{häufige Fehler} & \textcolor{gruen}{keiner}\\
\code{go test} & Prüfungen, Abdeckung & \textcolor{gruen}{keiner}\\
\code{-race} & \textbf{Wettläufe} & \textcolor{mittel}{Laufzeit}\\
\code{pprof} & \textbf{Profil} & \textcolor{gruen}{gering}\\
\code{go doc} & Dokumentation & \textcolor{gruen}{keiner}\\
\bottomrule
\end{tabular}}
\end{center}

\begin{center}
\begin{minipage}{0.88\textwidth}
\kursiv{Sechs Werkzeuge, keine Installation, keine Konfiguration --- sie sind Teil der Sprachauslieferung. Das ist der praktische Unterschied zu Band 21, wo dieselben Aufgaben eine Werkzeugkette verlangen, die jemand einrichten muss.}
\end{minipage}
\end{center}

\textbf{Was das für die Verbreitung bedeutet.}

\begin{center}
\renewcommand{\arraystretch}{1.45}
\begin{tabular}{@{}P{3.0cm}P{2.2cm}P{2.4cm}P{1.6cm}@{}}
\toprule
\textbf{Praxis} & \textbf{In C} & \textbf{In Go} & \\
\midrule
Formatierer läuft & $\approx35\,\%$ & \textcolor{gruen}{$\approx\mathbf{99\,\%}$} & \\
Statische Prüfung & $\approx25\,\%$ & \textcolor{gruen}{$\approx\mathbf{85\,\%}$} & \\
Wettlaufprüfer & $<10\,\%$ & \textcolor{gruen}{$\approx60\,\%$} & \\
\bottomrule
\end{tabular}
\end{center}

\begin{center}
\kursiv{Die Zahlen sind Schätzungen und die Größenordnung ist belastbar: Ein Werkzeug, das mitgeliefert wird und ohne Einrichtung läuft, wird benutzt. Eines, das man aufsetzen muss, oft nicht --- die Lehre aus Band 21, Kapitel~10.}
\end{center}

\begin{technik}
\boxtitel{Was \code{go vet} findet}
\begin{center}
{\small
\renewcommand{\arraystretch}{1.45}
\begin{tabular}{@{}P{3.2cm}P{2.4cm}P{2.4cm}@{}}
\toprule
\textbf{Fehler} & \textbf{Sonst bemerkt?} & \\
\midrule
Formatzeichen passt nicht & \textcolor{akzent}{zur Laufzeit} & \textcolor{gruen}{gefunden}\\
Struktur mit Sperre kopiert & \textcolor{akzent}{\textbf{nie}} & \textcolor{gruen}{\textbf{gefunden}}\\
Schleifenvariable in Goroutine & \textcolor{akzent}{selten} & \textcolor{gruen}{gefunden}\\
Fehler nicht geprüft & \textcolor{akzent}{nie} & \textcolor{gruen}{teils}\\
Unerreichbarer Code & \textcolor{akzent}{nie} & \textcolor{gruen}{gefunden}\\
\bottomrule
\end{tabular}}
\end{center}

\begin{center}
\begin{minipage}{0.86\textwidth}
\kursiv{Die zweite Zeile ist die wertvollste --- eine kopierte Sperre schützt nichts und fällt sonst nie auf. Die dritte war bis Go 1.22 eine der häufigsten Fehlerquellen überhaupt; seither ist sie durch eine Sprachänderung beseitigt.}
\end{minipage}
\end{center}
\end{technik}
\end{merksatz}

\begin{merksatz}
\boxtitel{Das Profil}
\begin{center}
\renewcommand{\arraystretch}{1.5}
\begin{tabular}{@{}P{2.6cm}P{2.6cm}P{3.0cm}@{}}
\toprule
\textbf{Art} & \textbf{Zeigt} & \\
\midrule
Rechenzeit & \textbf{wo die Zeit vergeht} & \\
Speicher & \textbf{wer belegt} & \textcolor{gruen}{Halde}\\
Blockierung & \textbf{wo gewartet wird} & \textcolor{gruen}{Kanäle, Sperren}\\
Goroutinen & \textbf{wo sie stehen} & \textcolor{gruen}{\textbf{Kap.~\ref{kap:schiefgeht}}}\\
Sperrenstreit & wer wartet auf wen & \\
\bottomrule
\end{tabular}
\end{center}

\begin{center}
\begin{minipage}{0.88\textwidth}
\kursiv{Die dritte und vierte Art gibt es in anderen Sprachen selten und sie sind für nebenläufige Programme die wichtigsten: Ein Programm, das nur wartet, verbraucht keine Rechenzeit --- ein Rechenzeitprofil zeigt dann nichts, obwohl das Problem offensichtlich ist.}
\end{minipage}
\end{center}

\textbf{Im Betrieb.}

\begin{center}
\begin{minipage}{0.9\textwidth}
\code{import \_ "net/http/pprof"}\\
\\
\code{go func() \{}\\
\code{\ \ \ \ log.Println(http.ListenAndServe("localhost:6060", nil))}\\
\code{\}()}
\end{minipage}
\end{center}

\begin{center}
\renewcommand{\arraystretch}{1.45}
\begin{tabular}{@{}P{3.0cm}P{2.2cm}P{2.8cm}@{}}
\toprule
\textbf{Eigenschaft} & \textbf{Wert} & \\
\midrule
Zeilen & \textcolor{gruen}{$\mathbf{4}$} & \\
Laufzeitkosten & \textcolor{gruen}{$\approx0$} & \textbf{bis abgefragt}\\
Im Betrieb nutzbar & \textcolor{gruen}{\textbf{ja}} & \textcolor{gruen}{Band 21, Kap.~12}\\
Sicherheit & \textcolor{akzent}{\textbf{nur lokal binden}} & \textcolor{akzent}{\textbf{nie nach außen}}\\
\bottomrule
\end{tabular}
\end{center}

\begin{center}
\kursiv{Vier Zeilen, und man kann ein laufendes Programm im Betrieb befragen --- das Mittel, das Band 21 vermisst. Die letzte Zeile ist keine Formalie: Die Profilschnittstelle gibt Speicherinhalte preis und gehört niemals ins offene Netz.}
\end{center}

\begin{warnung}
\boxtitelr{Der Formatierer ersetzt kein Urteil}
\begin{center}
\renewcommand{\arraystretch}{1.5}
\begin{tabular}{@{}P{3.0cm}P{2.2cm}P{2.8cm}@{}}
\toprule
\textbf{Werkzeug findet} & \textbf{Werkzeug findet nicht} & \\
\midrule
Formatierung & \textbf{schlechte Namen} & \\
Wettläufe & \textbf{falsches Datenmodell} & \textcolor{akzent}{Band 20, Kap.~2}\\
Kopierte Sperren & \textbf{zu große Funktionen} & \\
Ungenutzte Variablen & \textbf{fehlende Zerlegung} & \\
\bottomrule
\end{tabular}
\end{center}

\begin{center}
\begin{minipage}{0.88\textwidth}
\kursiv{Die rechte Spalte enthält alles, was Band 20 behandelt --- und kein Werkzeug findet davon etwas. Eine grüne Prüfstrecke sagt, dass die bekannten Fehlerklassen fehlen; sie sagt nichts darüber, ob der Entwurf trägt.}
\end{minipage}
\end{center}
\end{warnung}
\end{merksatz}

\begin{aufgabe}
\begin{enumerate}[label=(\alph*)]
\item Richten Sie die Prüfstrecke ein.
\item Finden Sie ein Leistungsproblem mit dem Profil.
\item Bewerten Sie zusätzliche Prüfer.
\end{enumerate}

\loesung
\textbf{(a) Die Prüfstrecke.}

\begin{center}
{\small
\renewcommand{\arraystretch}{1.45}
\begin{tabular}{@{}P{1.0cm}P{3.4cm}P{1.8cm}P{2.2cm}@{}}
\toprule
\textbf{St.} & \textbf{Schritt} & \textbf{Dauer} & \textbf{Wann}\\
\midrule
$1$ & \code{gofmt -l} --- meldet Abweichungen & $\mathbf{1}$ s & \textbf{jede Übergabe}\\
$2$ & \code{go vet} & $6$ s & jede Übergabe\\
$3$ & \code{go build} & $9$ s & jede Übergabe\\
$4$ & \code{go test} & $14$ s & jede Übergabe\\
$5$ & \textbf{\code{go test -race}} & $92$ s & \textbf{jede Übergabe}\\
$6$ & Zusatzprüfer & $40$ s & nächtlich\\
\bottomrule
\end{tabular}}
\end{center}

\begin{center}
\begin{minipage}{0.88\textwidth}
\kursiv{Die ersten fünf Stufen dauern zusammen gut zwei Minuten --- innerhalb der Grenze aus Band 17, Kapitel~10. Der Wettlaufprüfer ist der teuerste Posten und der wichtigste; wer ihn nach hinten schiebt, findet Wettläufe erst nachts.}
\end{minipage}
\end{center}

\begin{center}
\renewcommand{\arraystretch}{1.45}
\begin{tabular}{@{}P{3.0cm}P{2.2cm}P{2.8cm}@{}}
\toprule
\textbf{Einrichtung} & \textbf{Aufwand} & \\
\midrule
Einmalig & \textcolor{gruen}{$\approx\mathbf{2}$ \textbf{h}} & \textbf{Band 21: $16$ h}\\
Altlast\-meldungen & \textcolor{gruen}{$\approx30$} & \textcolor{gruen}{Band 21: $400$}\\
Konfiguration & \textcolor{gruen}{fast keine} & \\
\bottomrule
\end{tabular}
\end{center}

\begin{center}
\kursiv{Zwei Stunden gegen sechzehn, dreißig Altlastmeldungen gegen vierhundert. Der Unterschied ist nicht Sprachqualität, sondern die Entscheidung, die Werkzeuge mitzuliefern und den Formatierer von Anfang an verbindlich zu machen.}
\end{center}

\medskip
\textbf{(b) Ein Leistungsproblem finden.}

\textbf{Das Bild.} Die Übertragung schafft $180$ k/s statt der erwarteten $500$ k/s.

\begin{center}
{\small
\renewcommand{\arraystretch}{1.45}
\begin{tabular}{@{}P{1.0cm}P{3.4cm}P{2.0cm}P{2.0cm}@{}}
\toprule
& \textbf{Schritt} & \textbf{Dauer} & \textbf{Befund}\\
\midrule
$1$ & Rechenzeitprofil & $3$ min & \textcolor{akzent}{$8\,\%$ Last}\\
$2$ & \textbf{Blockierungsprofil} & $3$ min & \textcolor{gruen}{\textbf{Kanal}}\\
$3$ & Zeile ansehen & $5$ min & \textbf{Puffer $0$, ein Leser}\\
$4$ & Vier Leser starten & $10$ min & \textcolor{gruen}{$\mathbf{495}$ \textbf{k/s}}\\
\bottomrule
\end{tabular}}
\end{center}

\begin{center}
\begin{minipage}{0.88\textwidth}
\kursiv{Das Rechenzeitprofil zeigt acht Prozent Last --- also kein Rechenproblem. Erst das Blockierungsprofil zeigt, wo gewartet wird, und das ist bei einem nebenläufigen Programm fast immer die richtige Frage.}
\end{minipage}
\end{center}

\begin{center}
\kursiv{Einundzwanzig Minuten von der Meldung zur Behebung --- und der entscheidende Schritt war, das richtige Profil zu wählen. Wer nur Rechenzeit misst, sucht stundenlang an der falschen Stelle.}
\end{center}

\medskip
\textbf{(c) Zusätzliche Prüfer bewerten.}

\begin{center}
{\small
\renewcommand{\arraystretch}{1.45}
\begin{tabular}{@{}P{3.0cm}P{2.2cm}P{2.6cm}@{}}
\toprule
\textbf{Prüfer} & \textbf{Genauigkeit} & \textbf{Urteil}\\
\midrule
\code{go vet} & \textcolor{gruen}{$\approx95\,\%$} & \textcolor{gruen}{\textbf{immer}}\\
Wettlaufprüfer & \textcolor{gruen}{$\mathbf{100\,\%}$} & \textcolor{gruen}{\textbf{immer}}\\
Stilprüfer & \textcolor{mittel}{$\approx60\,\%$} & \textcolor{mittel}{auswählen}\\
Komplexitätsmaße & \textcolor{akzent}{$\approx30\,\%$} & \textcolor{akzent}{Hinweis, keine Regel}\\
Alle Prüfer eines Bündels & \textcolor{akzent}{$\approx35\,\%$} & \textcolor{akzent}{\textbf{Rauschen}}\\
\bottomrule
\end{tabular}}
\end{center}

\begin{center}
\begin{minipage}{0.88\textwidth}
\kursiv{Die letzte Zeile ist der häufigste Fehler beim Einrichten: Man schaltet alle vorhandenen Prüfer ein, bekommt tausend Meldungen und schaltet nach zwei Wochen alle wieder ab. Dieselbe Rechnung wie in Band 19, Kapitel~4 und Band 21, Kapitel~10.}
\end{minipage}
\end{center}

\begin{center}
\kursiv{Die beiden ersten Zeilen genügen für den Anfang. Weitere Prüfer nimmt man einzeln auf, wenn sie im eigenen Bestand echte Funde liefern --- gemessen an einer Stichprobe, nicht nach Empfehlung.}
\end{center}

\begin{warnung}
\boxtitelr{Der Wettlaufprüfer braucht Prüffälle}
\begin{center}
\renewcommand{\arraystretch}{1.5}
\begin{tabular}{@{}P{3.0cm}P{2.2cm}P{2.8cm}@{}}
\toprule
\textbf{Voraussetzung} & \textbf{Erfüllt?} & \\
\midrule
Nebenläufiger Pfad geprüft & \textcolor{mittel}{\textbf{oft nicht}} & \textcolor{akzent}{\textbf{die Lücke}}\\
Genug Durchläufe & \textcolor{mittel}{teils} & \\
Realistische Last & \textcolor{akzent}{selten} & \\
\bottomrule
\end{tabular}
\end{center}

\begin{center}
\begin{minipage}{0.88\textwidth}
\kursiv{Der Prüfer sieht nur, was ausgeführt wird --- dieselbe Einschränkung wie bei den Sanitizern in Band 21. Ein Wettlauf in einem Pfad, den kein Prüffall erreicht, bleibt unentdeckt, und die grüne Meldung wiegt in falscher Sicherheit.}
\end{minipage}
\end{center}

\begin{center}
\kursiv{Die Abhilfe ist eine Prüfung, die den nebenläufigen Pfad tatsächlich belastet: hundert Goroutinen, tausend Durchläufe, mit \code{t.Parallel}. Zwanzig Zeilen, und der Prüfer bekommt etwas zu sehen.}
\end{center}
\end{warnung}
\end{aufgabe}
\vspace{-4pt}
\noindent Die Werk\-zeug\-do\-ku\-men\-ta\-tion\zit{27} beschreibt das Profil und die Prüfer; der Wett\-lauf\-prü\-fer\zit{28} ist eigens do\-ku\-men\-tiert; \enquote{Effective Go}\zit{7} nennt die Ge\-wohn\-hei\-ten, die die Werk\-zeu\-ge voraussetzen.
